
en procédant par récurrence décroissante, supposons qu'elle est connue pour @@M0@@. Quitte à remplacer @@M1@@ par @@M2@@, on peut supposer que @@M3@@ est @@M4@@-plat pour @@M5@@. Comme @@M6@@ est plat sur @@M7@@, il résulte alors de EGA III 6.9.8 que @@M8@@ commute à tout changement de base. Par conséquent la représentabilité de @@M9@@ résulte du lemme suivant (voir : J.P. Murre, Representation of unramified functors, Séminaire Bourbaki~; Mai 1965, p.294-11, th.2) :

\textbf{Lemme 4.5.} \uline{Soient} @@M10@@ \uline{un morphisme propre de présentation finie,} @@M11@@ \uline{un faisceau de présentation finie sur} @@M12@@\uline{. Alors le sous-foncteur} @@M13@@ \uline{de} @@M14@@ \uline{qui rend} @@M15@@ \uline{plat sur} @@M16@@ \uline{est représentable par un monomorphisme surjectif de présentation finie (qui n'est bien sûr pas une immersion en général).}

Pour représenter la condition 3) de 4.3, on note que d'après 1) et 2), @@M17@@ et @@M18@@ commutent aux changements de base @@M19@@, et donc @@M20@@ y commute aussi~; ensuite on applique 4.5.

Pour représenter 4), on note que d'après 3) et le fait que @@M21@@ et @@M22@@ sont plats et commutent aux changements de base, @@M23@@ commute aussi~; ensuite on applique le

\textbf{Lemme 4.6.} \uline{Soient} @@M24@@ \uline{un morphisme propre de présentation finie,} @@M25@@ \uline{un homomorphisme de} @@M26@@\uline{-Modules de présentation finie, avec} @@M27@@ \uline{plat sur} @@M28@@\uline{. Alors le sous-foncteur} @@M29@@ \uline{de} @@M30@@ \uline{qui rend} @@M31@@ \uline{un isomorphisme est représentable par un sous-schéma ouvert de} @@M32@@\uline{, de présentation finie.}

On se ramène comme d'habitude au cas @@M33@@ noethérien, en utilisant EGA IV 8 et 11.2.6, ce qui nous dispense de nous préoccuper de la présentation finie. Représentons d'abord le foncteur qui rend @@M34@@ un épimorphisme,
